\section{Results}
\label{sec:results}

\paragraph{Reading the table.} All numbers below are on real
AffectNet~\cite{mollahosseini2017affectnet} valence--arousal under the protocol of
Section~\ref{sec:experiments}: three seeds for the baseline, Gaussian NLL, and
evidential models, a single seed for split conformal, with the corruption-based
shift as the out-of-distribution condition. We report the mean over seeds and the
population standard deviation. No synthetic data enters any reported result.

Table~\ref{tab:comparison} compares the four methods across accuracy,
calibration, selective prediction, and behaviour under shift.

\begin{table}[h]
  \centering
  \small
  \input{comparison_table.tex}
  \caption{Method comparison on AffectNet-VA. Mean$_{\pm\text{std}}$ over three
  seeds ($1337/7/21$); $^{\ddagger}$conformal is a single seed and is reported
  without error bars. CCC: mean concordance over valence and arousal (higher
  better). RMSE: mean over both axes (lower better). ECE: regression calibration
  error (lower better). Sharp.: mean predictive interval width, informative
  only alongside ECE and therefore never ranked. AURC: area under the risk--coverage
  curve (lower better). OOD-AUROC: separation of clean from corrupted inputs by
  the uncertainty score (higher better). $\Delta$ECE: calibration degradation
  from clean to corrupted inputs; all uncertainty-aware methods lie within seed
  noise of zero, so this column is likewise not ranked. $^{\dagger}$The
  deterministic baseline emits no uncertainty, so its ECE, sharpness, OOD-AUROC
  and $\Delta$ECE are degenerate rather than competitive and are excluded from
  ranking. \textbf{Bold} marks the best value: over all methods for the accuracy
  columns, and over the uncertainty-aware methods elsewhere.}
  \label{tab:comparison}
\end{table}

\paragraph{Uncertainty carries no accuracy penalty.} Evidential regression attains
the highest concordance (CCC $0.692$) while its RMSE ($0.286$) is within noise of
every other method ($\approx 0.28$--$0.29$). Adding a principled uncertainty head
therefore does not cost predictive accuracy. We deliberately do not claim an
accuracy \emph{gain}: the spread across methods ($0.657$--$0.692$ CCC) is
comparable to the spread across seeds within a method (up to $\pm0.017$), so on
this axis the four methods should be read as equivalent. That equivalence is the
precondition for the rest of our argument: uncertainty is obtained
\emph{for free} on the accuracy axis, so the question is purely how much
\emph{trustworthiness} each method adds, and there the methods are \emph{not}
equivalent at all.

\paragraph{Accuracy in context.} The claim that uncertainty is free on the
accuracy axis only carries weight if the underlying predictor is not weak, so we
position it against published AffectNet numbers. Our evidential model reaches
CCC $0.802_{\pm.002}$ on valence and $0.581_{\pm.015}$ on arousal. The AffectNet
authors' own CNN baseline reports $0.541$ and $0.450$ on their test set
($0.60$ and $0.34$ on their validation set)~\cite{mollahosseini2017affectnet},
and EmoNet~\cite{toisoul2021estimation}, a substantially larger purpose-built
architecture, reports $0.82$ and $0.75$ on a cleaned AffectNet test set. Our
valence accuracy is therefore close to the strongest published figure, while our
arousal accuracy falls between the dataset's own baseline and EmoNet, matching
the widely observed asymmetry that arousal is the harder axis to recover from a
single still image. We stress that these are \emph{indicative} rather than exact
comparisons: the public redistribution we train on defines its own
Train/Validation/Test partition rather than the authors' official split, and
EmoNet's figures are computed on a cleaned subset under a different protocol. The
point is narrow: the calibration comparison in Table~\ref{tab:comparison}
rests on a predictor of representative strength rather than a crippled one.

\paragraph{Calibration and selective prediction.} The deterministic baseline
provides no usable uncertainty: its expected calibration error is degenerate
($0.500$) and its sharpness is zero because it emits a point estimate. Among the
uncertainty-aware methods, evidential regression is the best calibrated
(ECE $0.011$, versus $0.024$ for Gaussian NLL and $0.015$ for conformal), a
better than twofold reduction over Gaussian NLL and the axis on which the
methods separate most clearly. It reaches that calibration without paying for it in
sharpness: its intervals ($0.245$) are as tight as any method's and tighter than
Gaussian NLL's ($0.270$), so it achieves tight intervals \emph{and} correct
coverage rather than trading one for the other. On selective prediction the three
uncertainty-aware methods are effectively tied (AURC $0.047$--$0.048$); all three
are far ahead of the baseline's $0.077$, which is the comparison that matters,
since it is the baseline a practitioner would otherwise deploy.
Figures~\ref{fig:rc} and~\ref{fig:rel} show the corresponding risk--coverage curve
and reliability diagram.

\begin{figure}[h]
  \centering
  \begin{minipage}{0.32\textwidth}
    \includegraphics[width=\linewidth]{evidential_risk_coverage.png}
    \subcaption{Risk--coverage}\label{fig:rc}
  \end{minipage}\hfill
  \begin{minipage}{0.32\textwidth}
    \includegraphics[width=\linewidth]{evidential_reliability.png}
    \subcaption{Reliability}\label{fig:rel}
  \end{minipage}\hfill
  \begin{minipage}{0.32\textwidth}
    \includegraphics[width=\linewidth]{evidential_va_scatter.png}
    \subcaption{VA, coloured by uncertainty}\label{fig:scatter}
  \end{minipage}
  \caption{Evidential model. (a) Risk falls monotonically as the most uncertain
  inputs are abstained on, falling roughly fourfold from $0.084$ at full coverage
  to $0.022$ at $5\%$, so the uncertainty score ranks errors usefully rather than
  merely being well calibrated on average. (b) Empirical coverage tracks the
  nominal level closely across the whole range, with mild overconfidence only in
  the extreme tail. (c) Uncertainty is highest in the \emph{interior} of the VA
  plane, where low-intensity and ambiguous expressions live, and lowest for the
  confident high-magnitude predictions on the perimeter. The model is least sure
  where human annotators also disagree most. The predictions also
  show the range compression typical of VA regression, spanning roughly
  $[-0.9, 0.8]$ in valence rather than the full $[-1,1]$.}
\end{figure}

\paragraph{Behaviour under shift.} The corruption shift is where the value of
calibrated uncertainty becomes concrete. All three uncertainty-aware methods
separate corrupted from clean inputs by their uncertainty score well above chance
and within noise of one another (OOD-AUROC $0.606$--$0.613$). The baseline, by
contrast, does not merely fail; it fails \emph{unpredictably}. Its three seeds
score $0.001$, $0.001$ and $0.987$, averaging to $0.329 \pm 0.465$. Lacking any
uncertainty head, its ``score'' is an artefact with no consistent sign, so a
practitioner could not tell from a single training run even which direction it
points. A point estimate cannot flag its own degradation. Equally important, the
uncertainty-aware models stay calibrated under shift: their ECE degrades by at
most $|\Delta\text{ECE}| \le 0.006$, so the abstention threshold calibrated on
clean data remains meaningful on corrupted data. In other words, the system
raises its hand when the input drifts rather than failing silently, which is the
behaviour a deployed affect estimator needs.

\paragraph{Summary.} Across accuracy, calibration, selective prediction, and shift
robustness, evidential regression is the strongest single-forward-pass method.
Its margin is not in accuracy or in selective prediction, where it is tied with
the alternatives, but in calibration, the quantity from which the abstention
threshold is actually derived. Gaussian NLL is a close and cheaper aleatoric-only
alternative with equal OOD separation and roughly twice the calibration error, and
conformal prediction adds distribution-free coverage guarantees on top of the same
predictor. The deterministic baseline, competitive on raw accuracy, is unusable
wherever trustworthiness matters.

\paragraph{Limitations.} Conformal prediction is currently reported for a single
seed; extending it to the same three-seed protocol is immediate future work. We
also provide Monte-Carlo dropout and deep ensembles in the released harness but
exclude them from the main comparison, as MC-dropout requires a model trained with
dropout enabled rather than a reused evidential checkpoint. Finally, the shift
evaluation here is corruption-based; the cross-dataset
AffectNet$\rightarrow$AFEW-VA shift runs identically through the same harness once
that (registration-gated) data is in place.
